\section*{الفصل \ref{ch:g7:atoms-and-molecules} --- الذرات والجزيئات}

\begin{solution}{exo:g7:atoms-and-molecules:1}
\omterm{def:g7:atoms-and-molecules:atom}{الذرة} واحد من الجسيمات البالغة الصغر التي بُنيت منها كل المادة (ويوجد منها
نحو مئة نوع). \omterm{def:g7:atoms-and-molecules:molecule}{والجزيء} مجموعة من \omterm{def:g7:atoms-and-molecules:atom}{الذرات} مرتبط بعضها ببعض ارتباطًا وثيقًا، وهو
أصغر قطعة من مادة تبقى هي تلك
المادة.
\end{solution}

\begin{solution}{exo:g7:atoms-and-molecules:2}
\ce{H}، \ce{C}، \ce{N}، \ce{O}، \ce{Na}، \ce{Fe}.
\end{solution}

\begin{solution}{exo:g7:atoms-and-molecules:3}
\omterm{def:g7:atoms-and-molecules:atom}{ذرة} كربون واحدة وذرتا أكسجين: أي 3 \omterm{def:g7:atoms-and-molecules:atom}{ذرات} في المجموع.
\end{solution}

\begin{solution}{exo:g7:atoms-and-molecules:4}
الماء، والأكسجين، والميثان، والنيتروجين.
\end{solution}

\begin{solution}{exo:g7:atoms-and-molecules:5}
\ce{NO2}.
\end{solution}

\begin{solution}{exo:g7:atoms-and-molecules:6}
\ce{O2} \omterm{def:g7:atoms-and-molecules:molecule}{جزيء} واحد مكوّن من ذرتي أكسجين مرتبطتين؛ أما \ce{2O} فذرتا
أكسجين منفصلتان. و\ce{CO} \omterm{def:g7:atoms-and-molecules:molecule}{جزيء} من \omterm{def:g7:atoms-and-molecules:atom}{ذرة} كربون واحدة \omterm{def:g7:atoms-and-molecules:atom}{وذرة} أكسجين
واحدة (حرفان كبيران، أي رمزان)؛ أما \ce{Co} \omterm{def:g7:atoms-and-molecules:atom}{فذرة} واحدة من
الكوبالت (\omterm{def:g7:atoms-and-molecules:symbol}{رمز} واحد، حرف كبير يليه حرف صغير).
\end{solution}

\begin{solution}{exo:g7:atoms-and-molecules:7}
\ce{3H2O}: $3\times 2 = 6$ \omterm{def:g7:atoms-and-molecules:atom}{ذرات} هيدروجين و$3\times 1 = 3$ \omterm{def:g7:atoms-and-molecules:atom}{ذرات}
أكسجين. \ce{5O2}: لا هيدروجين، و$5\times 2 = 10$ \omterm{def:g7:atoms-and-molecules:atom}{ذرات} أكسجين.
\end{solution}

\begin{solution}{exo:g7:atoms-and-molecules:8}
الأسود كربون، والأحمر أكسجين: كربون واحد وأكسجينان، \ce{CO2}،
ثنائي أكسيد الكربون.
\end{solution}

\begin{solution}{exo:g7:atoms-and-molecules:9}
كان ماء دالتون \ce{HO}. \omterm{def:g7:atoms-and-molecules:molecule}{والجزيء} الحقيقي، \ce{H2O}، فيه \omterm{def:g7:atoms-and-molecules:atom}{ذرة} هيدروجين
إضافية: \omterm{def:g7:atoms-and-molecules:molecule}{فجزيء} دالتون كانت تنقصه \omterm{def:g7:atoms-and-molecules:atom}{ذرة} واحدة.
\end{solution}

\begin{solution}{exo:g7:atoms-and-molecules:10}
الماء النقي ليس خليطًا: فجزيئاته كلها \omterm{def:g7:atoms-and-molecules:molecule}{جزيئات} \ce{H2O}
متماثلة. \omterm{def:g4:air-a-mixture-of-gases:air}{والهواء} \omterm{def:g6:pure-substances-and-mixtures:mixture}{خليط}: فهو يحوي \omterm{def:g7:atoms-and-molecules:molecule}{جزيئات} نيتروجين، \omterm{def:g7:atoms-and-molecules:molecule}{وجزيئات}
أكسجين، \omterm{def:g7:atoms-and-molecules:atom}{وذرات} أرغون، وبضعة \omterm{def:g7:atoms-and-molecules:molecule}{جزيئات} من ثنائي أكسيد الكربون.
\end{solution}

\begin{solution}{exo:g7:atoms-and-molecules:11}
عشرة ملايين \omterm{def:g7:atoms-and-molecules:atom}{ذرة} تصنع \qty{1}{mm}، إذن \qty{1}{cm} $= \qty{10}{mm}$
تتطلب $10 \times 10\,000\,000 = 100\,000\,000$ \omterm{def:g7:atoms-and-molecules:atom}{ذرة}، أي مئة مليون \omterm{def:g7:atoms-and-molecules:atom}{ذرة}.
و$\qty{100}{m} = \num{100000}$~mm تتطلب $\num{100000}\times 10\,000\,000$
\omterm{def:g7:atoms-and-molecules:atom}{ذرة}، أي مليون مليون \omterm{def:g7:atoms-and-molecules:atom}{ذرة}.
\end{solution}

\begin{solution}{exo:g7:atoms-and-molecules:12}
$6 + 12 + 6 = 24$ \omterm{def:g7:atoms-and-molecules:atom}{ذرة}. الهيدروجين: $\frac{12}{24} = \qty{50}{\percent}$.
الكربون: $\frac{6}{24} = \qty{25}{\percent}$ (والأكسجين
\qty{25}{\percent} الباقية).
\end{solution}

\begin{solution}{pb:g7:atoms-and-molecules:1}
\textbf{1.} النيتروجين \ce{N2}، ذرتان؛ الأكسجين \ce{O2}، ذرتان؛ ثنائي
أكسيد الكربون \ce{CO2}، 3 \omterm{def:g7:atoms-and-molecules:atom}{ذرات}.

\textbf{2.} \omterm{def:g7:atoms-and-molecules:atom}{ذرات} الأرغون لا تتحد في \omterm{def:g7:atoms-and-molecules:molecule}{جزيئات}: فالأرغون مكوّن من \omterm{def:g7:atoms-and-molecules:atom}{ذرات}
منفردة.

\textbf{3.} النيتروجين: $78 \times 2 = 156$ \omterm{def:g7:atoms-and-molecules:atom}{ذرة}. الأكسجين: $21\times 2 =
42$ \omterm{def:g7:atoms-and-molecules:atom}{ذرة}.

\textbf{4.} $156 + 42 + 1 = 199$ \omterm{def:g7:atoms-and-molecules:atom}{ذرة}.

\textbf{5.} $\frac{156}{199} \approx 0.78$، أي نحو \qty{78}{\percent}.

\textbf{6.} في \omterm{def:g7:atoms-and-molecules:molecule}{جزيئات} الأكسجين: $16\times 2 = 32$؛ وفي \omterm{def:g7:atoms-and-molecules:molecule}{جزيئات} ثنائي أكسيد
الكربون: $4 \times 2 = 8$؛ وفي المجموع $32 + 8 = 40$.

\textbf{7.} ينقص $42 - 40 = 2$ من \omterm{def:g7:atoms-and-molecules:atom}{ذرات} الأكسجين: أي ما يعادل \omterm{def:g7:atoms-and-molecules:molecule}{جزيء} أكسجين
واحدًا. وقد ذهبتا إلى \omterm{def:g7:atoms-and-molecules:molecule}{جزيئات} الماء، \ce{H2O}، التي يصنعها الجسم من
هيدروجين غذائه؛ ويخرج هذا الماء جزئيًا بخارَ ماء في النَّفَس (وقد نُزع من
هذه الأرقام).

\textbf{8.} في الشهيق: 0. وفي الزفير: 4 (واحدة في كل \ce{CO2}).

\textbf{9.} من الغذاء الذي يستعمله الجسم: فالسكريات والدهون تحوي \omterm{def:g7:atoms-and-molecules:atom}{ذرات}
كربون، ويطلقها الجسم في ثنائي أكسيد الكربون الذي يزفره.

\textbf{10.} اختفت $21 - 16 = 5$ \omterm{def:g7:atoms-and-molecules:molecule}{جزيئات} أكسجين وظهرت 4 \omterm{def:g7:atoms-and-molecules:molecule}{جزيئات} من ثنائي
أكسيد الكربون.

\textbf{11.} النيتروجين: كرتان زرقاوان تربطهما ثلاث عصيّ. الأكسجين:
كرتان حمراوان تربطهما عصوان. ثنائي أكسيد الكربون: كرة سوداء بين كرتين
حمراوين، ترتبط بكل منهما بعصوين، على خط مستقيم. الأرغون: كرة واحدة زرقاء
فاتحة.

\textbf{12.} $156 + 32 + 12 + 1 = 201$ \omterm{def:g7:atoms-and-molecules:atom}{ذرة} (\omterm{def:g7:atoms-and-molecules:molecule}{جزيئات} ثنائي أكسيد الكربون
الأربعة تحوي 12 \omterm{def:g7:atoms-and-molecules:atom}{ذرة}). أي $201 - 199 = 2$ من \omterm{def:g7:atoms-and-molecules:atom}{الذرات} زيادةً على \omterm{def:g4:air-a-mixture-of-gases:air}{هواء}
الشهيق: 4 \omterm{def:g7:atoms-and-molecules:atom}{ذرات} كربون اكتُسبت من الغذاء، ناقص ذرتي الأكسجين اللتين خرجتا
في الماء.
\end{solution}
